% generated by analysis/r2_2_awarding_countries.py -- do not edit by hand
\begin{table}[H]
\centering \footnotesize \setlength{\tabcolsep}{3pt}
\caption{Awarding Countries of the 99 Prizes, Home Bias, and Where Laureates Are Based}
\label{tab:r2_2_awarding_countries}
\begin{tabular}{lrrrrrrrrr}
\toprule
Country & \begin{tabular}[c]{@{}c@{}}Prizes \\ given by \\ country\end{tabular} & \multicolumn{2}{c}{\begin{tabular}[c]{@{}c@{}}Recognitions awarded \\ by the country's \\ organizations\end{tabular}} & \multicolumn{3}{c}{\begin{tabular}[c]{@{}c@{}}Share of the country's \\ recognitions given to \\ scientists based there (\%)\end{tabular}} & \multicolumn{3}{c}{\begin{tabular}[c]{@{}c@{}}Share of all recognitions \\ received by scientists \\ based in the country (\%)\end{tabular}} \\
\cmidrule(lr){3-4}\cmidrule(lr){5-7}\cmidrule(lr){8-10}
 & & N & \% of all & Actual & \begin{tabular}[c]{@{}c@{}}Expected \\ from pool\end{tabular} & Ratio & Actual & \begin{tabular}[c]{@{}c@{}}Share \\ of pool\end{tabular} & Ratio \\
\midrule
United States & 33 & 485 & 33.2 & 76.9 & 65.1 & 1.2 & 65.0 & 65.0 & 1.0 \\
Sweden & 13 & 162 & 11.1 & 3.2 & 1.9 & 1.7 & 2.0 & 1.9 & 1.1 \\
Spain & 5 & 109 & 7.5 & 2.8 & 1.3 & 2.2 & 0.6 & 1.2 & 0.5 \\
Israel & 6 & 93 & 6.4 & 7.5 & 1.5 & 4.9 & 2.2 & 1.5 & 1.4 \\
Italy & 3 & 90 & 6.2 & 9.5 & 1.7 & 5.5 & 1.8 & 1.9 & 0.9 \\
United Kingdom & 7 & 87 & 6.0 & 42.9 & 12.0 & 3.6 & 13.6 & 13.0 & 1.0 \\
Japan & 5 & 85 & 5.8 & 22.6 & 3.0 & 7.5 & 4.4 & 3.4 & 1.3 \\
Canada & 2 & 62 & 4.2 & 11.3 & 3.3 & 3.4 & 4.6 & 4.3 & 1.1 \\
Germany & 5 & 56 & 3.8 & 30.9 & 8.4 & 3.7 & 10.0 & 8.1 & 1.2 \\
Norway & 4 & 55 & 3.8 & 0.0 & 0.4 & 0.0 & 0.1 & 0.4 & 0.3 \\
Hong Kong & 3 & 50 & 3.4 & 2.0 & 2.6 & 0.8 & 0.4 & 0.7 & 0.6 \\
Denmark & 1 & 33 & 2.3 & 9.1 & 0.9 & 10.3 & 1.5 & 1.2 & 1.3 \\
International & 3 & 26 & 1.8 & -- & -- & -- & -- & -- & -- \\
Netherlands & 4 & 16 & 1.1 & 6.2 & 3.0 & 2.1 & 1.7 & 2.8 & 0.6 \\
Saudi Arabia & 1 & 16 & 1.1 & 0.0 & 0.5 & 0.0 & 0.4 & 0.6 & 0.7 \\
Taiwan & 1 & 15 & 1.0 & 0.0 & 2.7 & 0.0 & 0.0 & 0.3 & 0.0 \\
Switzerland & 1 & 8 & 0.5 & 14.3 & 5.7 & 2.5 & 3.4 & 3.3 & 1.0 \\
Belgium & 1 & 7 & 0.5 & 0.0 & 1.1 & 0.0 & 1.1 & 1.3 & 0.9 \\
Finland & 1 & 6 & 0.4 & 16.7 & 0.9 & 17.9 & 0.2 & 0.6 & 0.3 \\
\midrule
\multicolumn{10}{l}{\emph{Countries awarding none of the 99 prizes (the 5 with the largest pool shares), and all others:}} \\
China (mainland) & 0 & 0 & 0.0 & -- & -- & -- & 1.5 & 3.5 & 0.4 \\
Australia & 0 & 0 & 0.0 & -- & -- & -- & 2.3 & 3.5 & 0.7 \\
France & 0 & 0 & 0.0 & -- & -- & -- & 5.2 & 3.2 & 1.6 \\
Russia & 0 & 0 & 0.0 & -- & -- & -- & 1.3 & 1.0 & 1.3 \\
Austria & 0 & 0 & 0.0 & -- & -- & -- & 1.7 & 0.9 & 1.9 \\
All other countries (99) & 0 & 0 & 0.0 & -- & -- & -- & 6.8 & 6.2 & 1.1 \\
\bottomrule
\end{tabular}
\par\smallskip\begin{minipage}{\textwidth}\footnotesize\textit{Note}: This table counts prizes and 2015--2024 recognitions by the country of the awarding organization; the three international prizes (Fields Medal, ICTP Ramanujan Prize, IABSE Award of Merit) have no national home. Home shares cover the 1{,}399 recognitions of national prizes with a placeable laureate: actual is the share of laureates based in the awarding country by the award year; expected is the share of the field's benchmark pool (the 1{,}000 researchers with the most citations to small-team papers, per field and award year) based there. The last three columns cover all 1{,}424 placeable recognitions: the share received by laureates based in the country and the country's share of the benchmark pool, weighted by the recognitions' field-year mix; a researcher based in two countries counts in both. Ratios divide actual by expected or pool shares. Hong Kong, Macau and Taiwan are folded into Greater China in the home-bias columns (1.8\% of laureates against 4.2\% of the pool) but not in the received-share columns. The last block lists the 5 non-awarding countries with the largest pool shares; the remaining countries are pooled in one row.\end{minipage}
\end{table}
